\documentclass[a4paper, 12pt]{article}

\usepackage[english,russian]{babel}
\usepackage[T2A]{fontenc}


% Абзацы и геометрия страницы
\usepackage{indentfirst}
\usepackage{geometry}

% Математика
\usepackage{amsmath}
\usepackage{amsthm}
\usepackage{amsfonts}
\usepackage{amssymb}

% Графика
\usepackage{tikz}
\theoremstyle{definition}
\newtheorem{definition}{Определение}
\newtheorem{example}{Пример}
\newtheorem{remark}{Замечание}

\theoremstyle{plain}
\newtheorem{theorem}{Теорема}

\newcommand{\N}{\mathbb{N}}
\newcommand{\R}{\mathbb{R}}
\newcommand{\eps}{\varepsilon}
\newtheorem{axiom}{Аксиома}

\usepackage{cancel}

% Списки и таблицы
\usepackage{enumitem}
\usepackage{tabularx}

% Дополнительные пакеты
\usepackage{blindtext}
\usepackage{titlesec}
\usepackage{needspace}

% Ссылки
\usepackage[
    colorlinks=true,
    linkcolor=blue,
    urlcolor=blue,
    citecolor=blue
]{hyperref}

% Колонтитулы
\usepackage{fancyhdr}
\pagestyle{fancy}
\fancyhf{}

% Берём название из \section
\renewcommand{\sectionmark}[1]{\markright{#1}}

\fancyhead[L]{\rightmark}
\fancyhead[R]{\hyperref[contents]{Содержание}}

\fancyfoot[C]{\thepage}

\renewcommand{\headrulewidth}{0.4pt}


% Своя стрелка
\newcommand{\myarrow}{%
    \begin{tikzpicture}[baseline=-0.5ex, scale=0.35]
        \draw[->, line width=0.8pt]
            (0,0.5)
            .. controls (-0.2,0.2) and (-0.2,-0.2) ..
            (0.2,-0.2)
            .. controls (0.5,-0.2) and (0.8,-0.2) ..
            (1.2,-0.2);
    \end{tikzpicture}%
}

\setlength{\tabcolsep}{4pt}
\pagestyle{empty}

% Дата лекции справа, как в тетради
\newcommand{\lecdate}[1]{\vspace{-4pt}\hfill{\small\textit{#1}}\par}

\begin{document}
\begin{center}
{\large КОНСПЕКТ ЛЕКЦИЙ}

ПО ДИСЦИПЛИНЕ <<ИСТОРИЯ>>

\tableofcontents
\label{contents}

\pagebreak
\end{center}
\vspace{10pt}

% =====================================================================
\section{Древний мир}
\lecdate{17.09.2026}

Первые государства возникли в IV тыс. до н.\,э. на Древнем Востоке~---
цивилизации Древнего Египта и Месопотамии. Главная функция государства~---
\underline{экономическая}.

Такие государства, как Хеттское царство, Урарту, Ассирия, были созданы
кочевниками. Форма государства~--- \textbf{деспотия}. <<Мирные>> аграрные
деспотии регулярно совершали завоевательные походы.

Черты деспотии:
\begin{enumerate}[noitemsep]
  \item обожествление правителя;
  \item большая роль жречества;
  \item появление документально закреплённых законов;
  \item патриархальное рабство.
\end{enumerate}

\textbf{Иерархия}~--- последовательное расположение общественных слоёв или
служебных званий.

Основная масса населения~--- крестьяне.

\textbf{Патриархальное рабство}: рабы~--- младшие члены семьи. Труд рабов
вспомогательный, их жизнь регулировалась законом.

\subsection{Финикия}
\begin{enumerate}[noitemsep]
  \item кораблестроение;
  \item мореплавание;
  \item торговля.
\end{enumerate}

Причины колонизации:
\begin{enumerate}[noitemsep]
  \item нехватка земель;
  \item стремление купцов закрепиться на торговых маршрутах;
  \item попытка сгладить социальные отношения.
\end{enumerate}

\subsection{Древняя Индия}
Особенности Древней Индии:
\begin{enumerate}[noitemsep]
  \item специфика географического положения;
  \item деление на касты.
\end{enumerate}

\textbf{Каста}~--- группа людей, в которой человек находится до своей смерти
(брахманы, кшатрии, вайшьи, шудры, далиты).

Буддизм появился в Индии, но не стал религией государства.

\subsection{Древний Китай}
Особенности Древнего Китая:
\begin{itemize}[noitemsep]
  \item получила распространение практика сдачи экзаменов;
  \item \textbf{конфуцианство}~--- особая морально-этическая система,
  созданная Конфуцием;
  \item <<благородный муж>>~--- идеал чиновника;
  \item чиновники противопоставлялись знати;
  \item государство определяло цены на продукцию.
\end{itemize}

Цивилизации Востока уступают античной цивилизации.

\subsection{Античность (Греция)}
Неудобное положение: из-за горной местности трудно заниматься сельским
хозяйством. Растут оливки и виноград.

\textbf{Полис}~--- небольшое государство, в территорию которого входил город.

\textbf{Классическое рабовладение}: рабы считались <<говорящими орудиями>>.

\textbf{Аристократия}~--- власть знати (аристократов).

\textbf{Демократия}~--- выборы на конкурентной основе.

\textbf{Колонизация}~--- основание новых поселений, ставших независимыми
полисами. Развитие мореплавания и торговли.

\textbf{Греко-персидские войны}~--- борьба против персидского завоевания.

\subsection{Древний Рим}
Территория лучше подходит для занятия сельским хозяйством.

Устанавливается \textbf{республика}:
\begin{enumerate}[noitemsep]
  \item консулы;
  \item народные трибуны;
  \item сенат.
\end{enumerate}

\textbf{Нобилитет}~--- самые богатые граждане.

Социальное устройство:
\begin{enumerate}[noitemsep]
  \item противопоставление свободных рабам;
  \item отсутствие деления на граждан и неграждан.
\end{enumerate}

Кризис:
\begin{enumerate}[noitemsep]
  \item вторжение варваров;
  \item распространение христианства;
  \item кризис рабовладельческого хозяйства.
\end{enumerate}

\textbf{Варвары}~--- народы, жившие догосударственным строем (германцы,
кельты, славяне).

\textbf{Военная демократия}~--- обозначение власти на стадии перехода
к государству.

Великое переселение народов. Рим разделился на две части.

% =====================================================================
\section{Средние века}
\lecdate{24.09.2026}

\subsection{Раннее Средневековье}
Варварские королевства. <<Тёмные века>>. Упадок культуры. Происходит
отставание европейской цивилизации.

\textbf{Климатический минимум}: температура была на 1--1,5 градуса ниже
нынешней.

Уровень социального расслоения был невысок. Хирдманы и хускарлы~---
дружинники того времени.

Постепенно складываются феодальные отношения: крестьяне попадают в
зависимость от землевладельцев, а знать связывается между собой
вассальными отношениями.

\textbf{Оммаж} (клятва верности) и \textbf{инвеститура} (введение в
должность) становятся неотъемлемыми в социальных отношениях.

\textbf{Соседская община} (марка) приводит к ещё большему расслоению.
\textbf{Аллод}~--- отчуждаемая земельная собственность отдельных малых семей.
Ценность земли закрепляет варварская аристократия.

После смерти короля все его потомки должны были получить свою долю:
государственные земли считались личным имуществом монарха. Это приводило
к братоубийственным войнам.

<<\textbf{Салическая правда}>>~--- письменный закон, который был
зафиксирован на вульгарной латыни. Утратила актуальность к X веку.

Провозглашается \textbf{империя франков} со столицей в Ахене. Карла Великого
короновали в Риме. Держава Карла Великого занимала очень большие территории.

Основа войска~--- \textbf{рыцари} (профессиональные тяжеловооружённые конные
воины). За военную службу рыцари получали \textbf{бенефиций}~--- условную
земельную собственность. Рыцарская небольшая армия была практически
непобедима.

Начали строить \textbf{замки} (укреплённые жилища). Опираясь на них, можно
было обороняться от превосходящих сил.

\subsection{Византия}
Стабильная экономика, многочисленное городское население, развитие ремесла
и торговли. Были развиты торговые связи с Европой и Востоком.

Византийцы продолжали традиции Римской империи. Власть была в руках
\textbf{василевса}, но был ещё и \textbf{синклит}~--- совет высших чиновников.

Особенностью являлось сохранение сильной централизованной системы. Расходы
империи были большими.

В условиях внешней угрозы началась феодализация. Страна была поделена на
военно-административные округа (фемы) во главе со стратегами.

\subsection{Исламский мир}
Особенности ислама:
\begin{itemize}[noitemsep]
  \item \textbf{шариат}~--- регламентация жизни;
  \item \textbf{джихад}~--- пропаганда религии;
  \item \textbf{халифат}~--- феодально-теократическое государство.
\end{itemize}

Было совершено очень много завоеваний на Ближнем и Среднем Востоке,
в Северной Африке и Испании.

Причины успеха:
\begin{enumerate}[noitemsep]
  \item истощение Византии и Ирана;
  \item многочисленность армии;
  \item роль религии.
\end{enumerate}

Немусульмане платили дополнительный налог.

\subsection{Высокое Средневековье}
Три сословия:
\begin{enumerate}[noitemsep]
  \item духовенство;
  \item крестьяне;
  \item рыцари.
\end{enumerate}

\textbf{Феодализм}~--- социально-экономическая формация, предшествующая
капитализму. В основе~--- поземельные и межличностные отношения.

\textbf{Барщина}~--- работа в хозяйстве господина, \textbf{оброк}~---
денежные выплаты.

\textbf{Феодальная лестница}: сеньор, вассал, феод, лен (короли, герцоги,
графы, бароны, рыцари).

Развиваются города, натуральное хозяйство переходит к товарному.

Особенности развития городов:
\begin{enumerate}[noitemsep]
  \item борьба городов за независимость от сеньоров. \textbf{Коммуна}~---
  статус свободного города;
  \item особые гражданские отношения. Формируется \textbf{буржуазия}.
\end{enumerate}

\textbf{Корпоративность}: цехи (ремесленники) и гильдии (объединения
купцов).

Появились пахотные земли. \textbf{Десятина}~--- $\frac{1}{10}$ доля дохода,
которая шла в церковь (налог).

\textbf{Индульгенция}~--- отпущение грехов.

Преступления, связанные с грехом, подлежали суду церкви:
\begin{enumerate}[noitemsep]
  \item ересь;
  \item вероотступничество;
  \item нарушение супружеской верности;
  \item лжесвидетельство.
\end{enumerate}

\textbf{Интердикт}~--- временное запрещение всех церковных действий.
\textbf{Святая инквизиция}~--- трибунал.

\subsection{Дальний Восток в Средние века}
Крестьян-общинников было очень много.

Особенности развития:
\begin{enumerate}[noitemsep]
  \item государство~--- верховный собственник всех земель;
  \item ни одна религия не закрепилась как единственная государственная.
\end{enumerate}

\textbf{Средневековая Индия:}
\begin{enumerate}[noitemsep]
  \item рост феодальных владений;
  \item кастовое деление общества;
  \item возникновение городов;
  \item из-за завоеваний~--- создание Делийского султаната (1206--1526);
  \item земли переходят к завоевателям.
\end{enumerate}

\textbf{Средневековый Китай:}
\begin{enumerate}[noitemsep]
  \item сельское хозяйство~--- основа экономики страны;
  \item сословное деление;
  \item господство трёх религий (конфуцианство, даосизм, буддизм).
\end{enumerate}

\textbf{Средневековая Япония:}
\begin{enumerate}[noitemsep]
  \item специфика географических условий;
  \item феодальные порядки. Сёгун~--- военный правитель, самураи~--- воины. \textbf{Сёгунат}~--- власть
  воинов.
\end{enumerate}

% =====================================================================
\section{Образование централизованных государств}
\lecdate{01.10.2026}

\begin{center}
\textbf{Высокое и позднее Средневековье}
\end{center}

Период феодальной раздробленности в Европе в XII--XIV вв. сменяется
формированием национальных государств, при этом во многих странах Западной
Европы сложилась новая \textbf{сословно-представительная монархия}. Наиболее
показательным в этом отношении является формирование централизованных
государств во Франции и Англии, поскольку Германия и Италия сохранили свою
раздробленность.

\subsection{Образование централизованного государства во Франции}
\begin{itemize}[noitemsep]
  \item В X в. Франция делилась на несколько крупных феодальных владений~---
  герцогства Нормандия, Аквитания, Бретань, Бургундия и графства Анжу,
  Шампань и др.
  \item В XII--XIV вв. происходит <<собирание>> французских земель
  королями~--- путём завоеваний, выгодных браков, конфискаций владений по
  обвинению в нарушении вассальной присяги.
\end{itemize}

\textbf{Филипп II Август} (1180--1223):
\begin{itemize}[noitemsep]
  \item увеличил королевский домен в 5 раз, присоединил Нормандию, Мэн, Анжу,
  часть Пуату и Тулузу;
  \item создаётся институт \textbf{прево} и \textbf{бальи}~--- королевских
  чиновников, управлявших его владениями за жалованье.
\end{itemize}

\textbf{Людовик IX Святой} (1226--1270):
\begin{itemize}[noitemsep]
  \item судебная реформа: создан институт профессиональных юристов~---
  легистов, появляется высшая судебная инстанция страны~--- Парижский
  парламент;
  \item денежная реформа: монеты королевской чеканки, в отличие от монет
  местного значения, имели хождение на территории всей Франции.
\end{itemize}

\textbf{Филипп IV Красивый} (1285--1314):
\begin{itemize}[noitemsep]
  \item присоединил к домену Шампань и Наварру;
  \item столкновение его с римским папой Бонифацием VIII по поводу обложения
  налогами церковных земель привело к \textbf{Авиньонскому пленению пап}
  (1309--1378).
\end{itemize}

В 1302 г. Филипп IV созывает \textbf{Генеральные штаты}~---
сословно-представительный орган власти, который включал в себя
представителей трёх сословий: духовенства, дворян и горожан.

\subsection{Образование централизованного государства в Англии}
После завоевания англами и саксами во время Великого переселения народов на
территории Британии образовались 7 королевств, которые объединились в одно
Английское королевство.

В 1066 г. нормандский герцог Вильгельм высадился в Британии с войском и,
разгромив англосаксонского короля Гарольда в битве при Гастингсе, был
коронован как король \textbf{Вильгельм Завоеватель}.

\textbf{Генрих II Плантагенет} (1154--1189):
\begin{itemize}[noitemsep]
  \item бароны были лишены права судить, были созданы королевский суд
  и институт присяжных.
\end{itemize}

В 1215 г. английские бароны, рыцари и горожане заставили короля Иоанна
Безземельного подписать \textbf{Великую хартию вольностей}.

В 1265 г. учреждается сословно-представительный орган власти~---
\textbf{Парламент} (в 1295 г. ставший постоянным).

Вопросы парламента: участие в издании законов, контроль над высшими
должностными лицами.

\subsection{Столетняя война (1337--1453 гг.)}
Это серия военных конфликтов между Англией и её союзниками и Францией и её
союзниками. Поводом к этим конфликтам были притязания на французский престол
Эдуарда III.
\begin{itemize}[noitemsep]
  \item В 1439 г. Карл VII вводит новый постоянный налог на землю
  и недвижимость~--- \textbf{талью}, взимаемый без согласия Генеральных штатов.
  \item С помощью нового налога финансируется постоянная армия.
\end{itemize}

\subsection{Абсолютизм}
\textbf{Абсолютизм}~--- близкая к диктатуре форма правления, когда монарху
принадлежит вся полнота неограниченной верховной власти.

Основные черты абсолютизма:
\begin{itemize}[noitemsep]
  \item создание разветвлённого бюрократического аппарата;
  \item уничтожение автономии городов;
  \item усиление карательных органов, появление постоянно действующей
  полиции.
\end{itemize}

\textbf{Война Алой и Белой розы} (1455--1485). Итогом войны явилось
физическое уничтожение почти всех знатных семей и победа дальнего
родственника Ланкастеров Генриха VII Тюдора, женившегося на Елизавете Йорк
и объединившего в гербе обе розы.

Своеобразием абсолютной монархии (в Англии) являлось:
\begin{itemize}[noitemsep]
  \item сохранение парламента;
  \item сохранение органов местного самоуправления.
\end{itemize}

\subsection{Полицентризм}
Примеры полицентризма:
\begin{itemize}[noitemsep]
  \item автономия итальянских городов являлась основой их процветания;
  \item княжества, входившие в Священную Римскую империю, были
  самостоятельными;
  \item к началу XII в. Германия была представлена множеством отдельных
  земель со столицами, которыми управляли герцоги, графы, князья.
\end{itemize}

В Средние века в Италии насчитывалось до 70 богатых городов, обладавших
значительной властью.

Внутри городов Италии шла постоянная борьба за власть между сторонниками
папы и императора. Иногда враждебными становились не только группы в
городе-коммуне, но и города-государства, поддерживающие ту или иную власть.
Часто в городах власть захватывали \textbf{кондотьеры}~--- командиры
наёмников, руководители военных отрядов.
\begin{itemize}[noitemsep]
  \item Разбогатевшие на торговле сукном и банковском деле синьоры стали
  силой захватывать власть в городах и устанавливать тиранию.
\end{itemize}

Во Флоренции к власти пришёл род аптекарей \textbf{Медичи}, заработавший
состояние на банковских сделках.

% =====================================================================
\section{Культура западноевропейского Средневековья}
\lecdate{08.10.2026}

\subsection{Общая характеристика западноевропейской средневековой культуры}
Культура Средневековья является неотъемлемым этапом мирового культурного
развития.

Условия культурного развития:
\begin{itemize}[noitemsep]
  \item феодальная форма собственности;
  \item господство натурального хозяйства;
  \item частые войны;
  \item эпидемии.
\end{itemize}

Достижения средневековой культуры:
\begin{itemize}[noitemsep]
  \item формирование современных европейских языков;
  \item расширение горизонта европейской культуры.
\end{itemize}

\begin{center}
\textbf{Мировоззренческие основы средневековой культуры}
\end{center}

Источники формирования культуры:
\begin{itemize}[noitemsep]
  \item наследие античности;
  \item культура варваров;
  \item традиции христианства.
\end{itemize}

\textbf{Наследие античности:}
\begin{itemize}[noitemsep]
  \item высокая правовая культура;
  \item христианство.
\end{itemize}

\textbf{Культура варваров:}
\begin{itemize}[noitemsep]
  \item языческое мировосприятие, чувство нерасчленённости природы, мира
  людей и мира богов.
\end{itemize}

\textbf{Христианство.} Появилось в I в.\footnote{\textit{Исправлено.} В
тетради~--- <<I в. до н.\,э.>>. Христианство возникло в I в. н.\,э.}
в восточных провинциях Римской империи.

Провозглашало: единобожие, равенство всех перед Богом, братство всех
народов. Христианство стало идеологическим стержнем духовной жизни
Средневековья и всей его культуры.

\textbf{Христианский Символ веры.} Христиане верят:
\begin{itemize}[noitemsep]
  \item в единого Бога Отца, Вседержителя, Творца неба и земли;
  \item в единую, святую, соборную и апостольскую церковь.
\end{itemize}

\textbf{Средневековая картина мира:}
\begin{itemize}[noitemsep]
  \item для христиан характерно противопоставление Бога и природы, неба
  и Земли. Мир виделся в соответствии с некоторой иерархической логикой, как
  симметричная схема, напоминающая две сложенные основаниями пирамиды;
  \item если Бог~--- олицетворение добра и любви, то Сатана~--- его
  противоположность.
\end{itemize}

Мирское время человека двигалось по замкнутому кругу.

\subsection{Сословные культуры Средневековья}
\begin{itemize}[noitemsep]
  \item высшее~--- духовенство;
  \item рыцарство~--- сословие воюющих;
  \item крестьянство~--- сословие работающих.
\end{itemize}

\textbf{1. Монастырская культура.} Монастыри~--- общины монахов.

Устав Василия Великого (свод правил): молитва, безбрачие, послушание
и повиновение.

Обитель Флавия Кассиодора: образованность~--- одна из первых христианских
добродетелей.

Борьба с еретиками объявляется высшим христианским подвигом.

Франциск Ассизский проповедовал святую бедность.

\textbf{2. Рыцарская культура.} Военно-аристократическое сословие, призванное
выполнять высшие государственные задачи.

\begin{center}
\textbf{Рыцарский кодекс чести}
\end{center}
\begin{itemize}[noitemsep]
  \item древность происхождения;
  \item верность сеньору;
  \item сила и храбрость.
\end{itemize}

\textbf{Внешний блеск}~--- особое внимание к оружию, атрибутике, символике
цвета.

\textbf{3. Народная культура.} Крестьянство составляло подавляющую часть
населения Европы. Высшим проявлением духа язычества были народные праздники,
в том числе и карнавалы.

\textbf{4. Городская культура.} С XI в. города становятся центрами духовной
жизни Западной Европы.

\subsection{Наука и образование}
Наука опиралась на абстрактное мышление. Основные направления:
физико-космологическое, наука о душе, алхимия.

\begin{center}
\textbf{Светское образование}
\end{center}

Кафедральные школы в крупнейших городах Европы превращались в университеты.
\textbf{Университет}~--- вуз.

Старейшие университеты Европы появились в XI в. в Италии.

\subsection{Литература и театр}
Неотъемлемой чертой городской культуры стали театральные представления.

\subsection{Архитектура и изобразительное искусство}
Включает в себя культуру <<варварских королевств>> и <<Каролингский
ренессанс>>.

Первый общеевропейский стиль~--- \textbf{романский стиль}. Основа стиля~---
храм-крепость и замок-крепость.

\textbf{Готический стиль} сменил романский. Ведущим архитектурным типом стал
собор. С исчезновением глухих стен исчезают и монументальные росписи. Их
заменяют \textbf{витражи}~--- живопись из кусков цветных стёкол.

% =====================================================================
\section{Культура эпохи Возрождения}
\lecdate{15.10.2026}

Это эпоха исключительного значения.

Основные черты: \textbf{гуманизм}, особое отношение к античности.

Периодизация Возрождения:
\begin{itemize}[noitemsep]
  \item Проторенессанс (XIII--XIV вв.);
  \item Раннее Возрождение (XV в.);
  \item Высокое Возрождение (конец XV~--- начало XVI в.).
\end{itemize}

\subsection{Новый этап развития науки}
Учения в области философии:
\begin{itemize}[noitemsep]
  \item стоицизм;
  \item неоплатонизм;
  \item пантеизм.
\end{itemize}

Другие достижения:
\begin{itemize}[noitemsep]
  \item Колумб в 1492 г. открыл Америку;
  \item Васко да Гама в 1498 г. открыл морской путь в Индию.
\end{itemize}

\subsection{Литература Возрождения}
Крупнейший представитель Проторенессанса~--- \textbf{Данте Алигьери}.

Истинным гуманистом был \textbf{Франческо Петрарка}.

\textbf{Франсуа Рабле}~--- писатель, врач, разносторонний учёный.

\subsection{Музыка}
В XVI в. в обиходе было много музыкальных инструментов: орган, клавесин,
виола, скрипка.

\subsection{Ренессансный театр}
В XVI в. в Италии родился первый профессиональный европейский театр.

\subsection{Изобразительное искусство эпохи Возрождения}
Для художников Возрождения изобразительным принципом стало подражание
природе: картина строилась как <<окно в мир>>, как нечто реально увиденное.

Система жанров:
\begin{itemize}[noitemsep]
  \item религиозно-мифологический;
  \item исторический;
  \item бытовой;
  \item пейзажный.
\end{itemize}

\begin{center}
\textbf{Раннее Возрождение}
\end{center}
\begin{itemize}[noitemsep]
  \item архитектор Филиппо Брунеллески;
  \item скульптор Донателло;
  \item живописец Мазаччо.
\end{itemize}

\textbf{Микеланджело Буонарроти} расписал потолок и создал фреску алтарной
стены <<Страшный суд>> Сикстинской капеллы в Ватикане. Он создал во
Флоренции капеллу Медичи~--- усыпальницу членов семьи Медичи, выполнил два
надгробия~--- Лоренцо и Джулиано.

\begin{center}
\textbf{Возрождение в Венеции}
\end{center}
Началось позже, длилось дольше, было своеобразнее.

% =====================================================================
\section{Культура Древней Руси}
\lecdate{22.10.2026}

\subsection{Материальная культура}
\textbf{Жилище славян} состояло из срубов, в княжеских и боярских жилищах
строились терема.

На лекции рассматривались иллюстрации: скотоводство у древних славян в
представлении художника, изделия кузнецов.

\textbf{Ювелирные техники}: скань, зернь, чернь, перегородчатая эмаль.

\subsection{Письменность и литература}
\textbf{Пергамент}~--- название происходит от города в Малой Азии~--- Пергама.

\textbf{Палимпсесты}~--- рукописи, написанные несколько раз на одном и том же
материале.

\textbf{Новгородский кодекс} (начало XI в.)~--- древнейшая книга Руси. Автор
неизвестен. \textbf{Цера}: состоит из липовых дощечек, покрытых воском,
с четырьмя страницами.

\textbf{Береста}~--- кора берёзы: менее прочный, но более дешёвый, чем
пергамент, материал.

\textbf{Украшение рукописей}: заставка, концовка, инициал, полевой цветок.

\textbf{Буквицы (инициалы)}~--- заглавная буква увеличенного размера, с
которой начинается текст.

\textbf{Литературные памятники}: фольклор, Остромирово Евангелие
(1056--1057).

\end{document}
